\documentclass[11pt]{article}
\usepackage[utf8]{inputenc}
\usepackage[T1]{fontenc}
\usepackage{amsmath,amssymb}
\usepackage{geometry}
\geometry{margin=1in}
\usepackage{hyperref}
\hypersetup{colorlinks=true,linkcolor=blue,urlcolor=blue}

\title{\textbf{VR-Forms --- v1.0.3 (revision note)}\\[2pt]
\large Two editorial corrections (VR v1.1.1 sync; ``no ontology'' scoped) and one added section: Things and Acts}
\author{Vitaly Reznik}
\date{2026 --- Version 1.0.3 (revision note)}

\begin{document}
\maketitle

\noindent\textbf{Scope.} This is the revision delta for VR-Forms v1.0.3. \emph{No construction,
definition, theorem, or axiom-profile of VR-Forms is altered.} Two editorial corrections and one
added expository section: (1) the recall of the parent work VR is brought up to date with VR v1.1.1;
(2) the ``no ontology'' language is scoped to ``no ontology of substance''; (3) a new section,
\emph{Things and Acts}, is added to the two-register discussion --- it names, and prices, what the
two registers already hold, resting on the cycle's existing axiom profiles (no new formal content).

\section*{(1) Correction: the parent-VR recall}

The introduction described VR as

\begin{quote}
``a minimal axiomatisation of arithmetic on \textbf{three primitives} $\{\emptyset,\to,t\}$ and
\textbf{four axioms}; equivalent to Peano arithmetic.''
\end{quote}

This is superseded. In VR v1.1.1 the propositional implication $\to$ (and the principles A1, A2 and
the two-valued logical register) are \textbf{removed}: $\to$ generated only a finite truth-function
algebra used by no VR theorem --- arithmetic, not logic --- and VR generates no logic of its own
(the logic it reasons with is metatheoretic, borrowed). The corrected recall reads

\begin{quote}
``a minimal \emph{axiom-free} reconstruction of arithmetic on \textbf{two primitives}
$\{\emptyset, t\}$ --- the nullary base $\emptyset$ and the unary succession $t$ --- with no axioms
of its own (membership is defined; the succession principle A3 is a theorem; induction A4 is the
recursor); equivalent to first-order Peano arithmetic
(Zenodo DOI 10.5281/zenodo.20688164, v1.1.1; concept 10.5281/zenodo.20212091).''
\end{quote}

\section*{(2) Correction: scoping the ``no ontology'' language}

Several places in VR-Forms state, without qualification, that ``the cycle ascribes no ontology''
(and the ``ontological register'' label is dropped on that basis). Read at face value this is a
\emph{blanket} denial of ontology, and it collides with the global membership relation $\in$ that
VR-Sets does carry. The intended and accurate claim is narrower: VR ascribes no ontology \emph{of
substance} --- no object that \emph{is} prior to operations; being is relocated into doing. It is not
the claim that the cycle ascribes no ontology \emph{at all}: the operational/set layer carries genuine
relational and structural commitments (global $\in$, potentialist infinity, relational identity),
which are \emph{tracked} (the two registers and the axiom tiers), not denied. Accordingly every
unqualified ``ascribes no ontology'' should read ``ascribes no ontology \emph{of substance}''; the
disclaimer, where it belongs, localises to the \emph{formal} register (``syntactic descriptions
without ontological commitment''), not to the cycle as a whole.

\section*{(3) New section: Things and Acts --- what inhabits the two registers}

\noindent\emph{To be added to the two-register discussion of the full paper.}

\paragraph{Acts and Things.}
An \emph{Act} is what an operation \emph{constructs}: an object obtained, and witnessed, by
performing operations --- a term, a witnessed bisimulation, a derivation that leans on nothing
posited. A \emph{Thing} is what an axiom \emph{posits}: an object, a totality, or an existence
principle introduced by \emph{postulation} rather than earned by construction (Hilbert's implicit
definition --- a Thing is whatever a consistent axiom makes it, relative to that axiom). The two
registers are exactly this contrast: the operational register is the register of Acts (built,
choice-free, becoming); the formal register is the register of Things (posited, axiomatic,
completed).

\paragraph{The one exception: the medium.}
Not everything unproven is a Thing. To state any axiom, or to derive any Act, one already reasons
--- with the connectives, the quantifiers, the rule of inference. This bare logic is not a Thing
but the \emph{medium}: it cannot be posited-or-not, for it is presupposed by every posit and every
derivation, and it is self-affirming --- to deny it one must already use it. Every genuine axiom,
by contrast, is \emph{optional} (a system may reject it and still reason) and is therefore a Thing.
In the cycle's ledger this puts the Act/Thing line at the empty profile: an Act depends on \emph{no}
axiom (\texttt{[]}); \texttt{propext} and \texttt{Quot.sound} are already Things, the cheapest ones;
\texttt{Classical.choice}, the axiom of infinity, the power set as a completed object, the aleph
cardinals and their consequences (e.g.\ Banach--Tarski) are the expensive ones. The operational
register is thus an Act built on the cheapest Things; the formal register adds the expensive ones.

\paragraph{How a Thing is admitted.}
A Thing is admitted not by construction but by a \emph{consistency check}: one exhibits a model ---
a structure in which the axiom holds (a theory is consistent iff it has a model) --- and thereby a
\emph{relative} tag $\mathrm{Con}(\text{base}) \to \mathrm{Con}(\text{base}+\text{Thing})$. The tag
is never absolute: by G\"odel's second theorem no consistent system strong enough for arithmetic
proves its own consistency. Hence the operational core, being axiom-free (\texttt{[]}), inherits the
host's consistency and adds no new burden, while a Thing is the first place a consistency question
arises at all.

\paragraph{Consistency is necessary but not sufficient; the cost discriminates.}
Every Thing that has a model is equally \emph{consistent}; consistency does not tell them apart.
What does is the \emph{kind of cost} the Thing carries. Three kinds recur across the cycle:
\begin{itemize}
  \item \textbf{strength} --- the axiom of infinity: completed $\omega$ proves $\mathrm{Con}(\mathrm{PA})$,
    so it cannot be justified from the finite base at all (G\"odel~II); a genuine leap in consistency
    strength;
  \item \textbf{impredicativity} --- the power set: to speak of ``all subsets of $\omega$'' one
    quantifies over the very totality being formed. Its profile is choice-free (\texttt{[propext]})
    and does \emph{not} reveal the circularity --- the impredicativity is a cost the machine profile
    fails to capture (finding~P-1);
  \item \textbf{non-constructivity} --- the axiom of choice: profile
    \texttt{[propext, Classical.choice, Quot.sound]} (the choice floor), an unexhibitable witness (no
    one displays a well-ordering of $\mathbb{R}$), and the non-measurable / Banach--Tarski pathology.
\end{itemize}
Thus $\wp(\omega)$ is a cheap, safe Thing and \textsc{ac} an expensive, pathological one, though both
are consistent: it is the cost, not consistency, that separates them --- and, as P-1 shows, not every
cost is visible to the axiom profile.

\paragraph{We track, we do not adopt.}
A Thing is placed in the formal register \emph{with} its tag and its cost; it is not adopted as an
operational truth. By the $O\!\to\!T\,/\,T\!\to\!O$ asymmetry a Thing does not transit down: the
completed power set, for instance, is used formally but has no operational correlate --- the
operational register carries only its \emph{becoming}, the non-enumerable $\wp$. To work with a
totality in VR is therefore not to build it but to posit it, exhibit its model, and pay --- in the
ledger --- for what it costs.

\section*{Why VR-Forms is otherwise unaffected}

VR-Forms (the operational two-register apparatus over VR-Sets, with the conservativity and transit
results) does not depend on the removed layer: its content references VR only through the operational
reading and the nullary base $\emptyset$, which is unchanged in v1.1.1. The Lean module
(\texttt{VRCycle.Forms}) contains no reference to the removed objects (\texttt{VRBool},
\texttt{impl}, A1, A2), and the build is unaffected. The ``propositional floor'' of the conservativity
theorem (\texttt{Forms/Conservativity.lean}) is an internal technical layer of Forms, unrelated to
VR's former $\to$ operator.

\vfill
\noindent\emph{Prepared with the assistance of Claude (Anthropic, Opus 4.8) in the cycle's Variant A
architecture; all mathematical content and positions are due to the author, Vitaly Reznik.}

\end{document}
